% Generated by scripts/build_theory_figures.py -- do not edit
\newcommand{\RevInfoNeed}{15}
\newcommand{\RevInfoLoadMin}{4603}
\newcommand{\RevInfoLoadMax}{6.9\times10^{6}}
\newcommand{\RevInfoLineMin}{4.9\times10^{6}}
\newcommand{\RevBusThirtyOneCross}{7.4}
\newcommand{\RevWorstOtherLoadPctTwenty}{1.9}
\newcommand{\RevBusThirtyOneMW}{9.2}
\newcommand{\RevResetEvents}{108}
\newcommand{\RevResetHarm}{84}
\newcommand{\RevResetHarmLoad}{33/54}
\newcommand{\RevResetHarmLine}{51/54}
\newcommand{\RevResetGainMin}{2.2\times10^{3}}
\newcommand{\RevResetFactorMin}{2}
\newcommand{\RevPredAlarmDelay}{0.20}
\newcommand{\RevObsLineDelay}{0.18}
\newcommand{\RevObsLoadDelay}{0.35}
\newcommand{\RevLoadExcessReports}{4}
\newcommand{\RevNoiseAgree}{96}
\newcommand{\RevNoiseLoadEvents}{12}
\newcommand{\RevNoiseLineEvents}{12}
\newcommand{\RevNoiseAddedRmsLow}{0.071}
\newcommand{\RevNoiseAddedRmsHigh}{0.141}
\newcommand{\RevNoiseOneLoad}{12/12}
\newcommand{\RevNoiseOneLine}{12/12}
\newcommand{\RevNoiseOneExtras}{0}
\newcommand{\RevNoiseOneNormalFA}{0}
\newcommand{\RevNoiseFiveLoad}{11/12}
\newcommand{\RevNoiseFiveLine}{12/12}
\newcommand{\RevNoiseFiveExtras}{0}
\newcommand{\RevNoiseFiveNormalFA}{0}
\newcommand{\RevNoiseTenLoad}{11/12}
\newcommand{\RevNoiseTenLine}{12/12}
\newcommand{\RevNoiseTenExtras}{0}
\newcommand{\RevNoiseTenNormalFA}{0}
\newcommand{\RevNoiseTwentyLoad}{12/12}
\newcommand{\RevNoiseTwentyLine}{12/12}
\newcommand{\RevNoiseTwentyExtras}{0}
\newcommand{\RevNoiseTwentyNormalFA}{0}
\newcommand{\RevNoiseTVEaLoad}{3/12}
\newcommand{\RevNoiseTVEaLine}{12/12}
\newcommand{\RevNoiseTVEaExtras}{0}
\newcommand{\RevNoiseTVEaNormalFA}{0}
\newcommand{\RevNoiseTVEbLoad}{0/12}
\newcommand{\RevNoiseTVEbLine}{10/12}
\newcommand{\RevNoiseTVEbExtras}{0}
\newcommand{\RevNoiseTVEbNormalFA}{0}
\newcommand{\RevNoiseNormalHours}{0.33}
\newcommand{\RevNomMsReport}{26.4}
\newcommand{\RevNomMsReportMax}{30.0}
\newcommand{\RevNomRecords}{6}
\newcommand{\RevAbstractNoise}{Uniform scaling of the calibrated noise preserved most attributions in the auxiliary stress test; independent per-PMU noise at \RevNoiseAddedRmsLow{}\% added RMS TVE degraded load attribution.}
\newcommand{\RevConclusionNoise}{The auxiliary stress test exposes sensitivity to cross-PMU noise structure. Estimating that covariance is a candidate extension, not a demonstrated remedy.}
\newcommand{\RevNoiseDiscussion}{Scaling the calibrated noise by up to twenty preserves all line attributions in this auxiliary cohort. One load intervention, a 4.1\% change at bus~26, is missed at $\times5$ and $\times10$ but not at $\times20$. Its amplitude exceeds the nominal reference at all three levels. That disagreement illustrates why a signature-energy crossing cannot guarantee a published attribution. The median load alarm delay changes from 0.35 to 0.36\,s. No scenario produces an extra alarm on event records or an alarm on the normal records.

Independent per-PMU noise has a larger effect: the bank attributes \RevNoiseTVEaLoad{} load interventions at $\sigma=0.0005$ and \RevNoiseTVEbLoad{} at $\sigma=0.001$, while the corresponding line counts are \RevNoiseTVEaLine{} and \RevNoiseTVEbLine{}. A matched-signature diagnostic using only the added independent noise still assigns substantial energy to these load responses (supplement). That restricted diagnostic excludes source search and the bank's covariance approximation; it does not establish achievable attribution accuracy. The bank fits one noise scale per signal family while retaining the reference cross-PMU correlation. Independent noise raises the fitted angle-noise scale and weakens load evidence under that model. Estimating cross-PMU covariance is therefore a candidate remedy that requires a separate controlled evaluation. These results do not establish readiness for field PMUs.}
